\section{对称性与约化}\label{sec:symmetry-and-reduction}

余切丛上的辛结构可以推广到任意流形.

\begin{RiemannDefinition}\label{def:symplectic-manifold}
\textbf{辛流形}是一个对 $(M,\omega)$, 其中 $M$ 是可微流形, $\omega\in\Omega^2(M)$ 非退化且闭.
\end{RiemannDefinition}

\begin{RiemannExample}\label{ex:orientable-surface-with-volume-form-is-symplectic}
若 $M$ 是可定向曲面, 且 $\omega\in\Omega^2(M)$ 是 $M$ 上的体积形式, 则 $(M,\omega)$ 为辛流形. 事实上, $\omega$ 必定非退化 (若存在非零 $v\in T_pM$ 使 $\iota(v)\omega=0$, 则 $\omega_p=0$), 并且显然有 $\mathrm{d}\omega=0$.
\end{RiemannExample}

第~\ref{sec:hamiltonian-mechanics}~节与第~\ref{sec:completely-integrable-systems}~节的所有定义及结果 (Hamilton 流及其性质, Liouville 定理与 Poincaré 回归定理, Poisson 括号, 完全可积系统和 Arnold–Liouville 定理) 都容易推广到任意辛流形. 事实上, 所有辛流形在局部都是一样的 (即在辛几何中没有对应于曲率的概念), 下面就来证明这一点.

\begin{RiemannTheorem}[Darboux]\label{thm:darboux-for-symplectic-manifolds}
设 $(M,\omega)$ 为辛流形, $p\in M$. 则存在 $p$ 附近的局部坐标 $(x^1,\ldots,x^n,\allowbreak p_1,\ldots,p_n)$, 使
\[
  \omega=\sum_{i=1}^n\mathrm{d}p_i\wedge\mathrm{d}x^i
\]
(特别地, $M$ 的维数必为偶数).
\end{RiemannTheorem}

\begin{RiemannProof}
先注意到 $\omega$ 具有上述形式, 当且仅当对 $i,j=1,\ldots,n$ 有 $\{x^i,x^j\}=\{p_i,p_j\}=0$ 和 $\{p_i,x^j\}=\delta_{ij}$ [参见 Exercise~\ref{ex:symmetry-and-reduction}(\ref{ex:canonical-brackets-and-symplectic-coordinates})].

显然必须有 $m:=\dim M\geq2$ (否则 $\omega=0$, 是退化的). 设 $P\in C^\infty(M)$ 为一个满足 $(\mathrm{d}P)_p\ne0$ 的函数, $X_P$ 为相应的 Hamilton 向量场, 并设 $T\subset M$ 为与 $(X_P)_p$ 不相切的一个超曲面 (参见 Fig.~\ref{fig:darboux-transverse-hypersurface}). 则在 $p$ 的某个邻域 $V$ 内, $X_P$ 与 $T$ 不相切. 必要时缩小 $V$, 可在 $V$ 上定义一光滑函数 $Q$, 使每个 $q\in V$ 都满足 $\psi_{-Q(q)}(q)\in T$, 其中 $\psi_t$ 是 $X_P$ 的流. 注意 $T\cap V=Q^{-1}(0)$, 这意味着 $X_Q$ 与 $T$ 相切, 因而 $\{(X_P)_p,(X_Q)_p\}$ 线性无关. 这说明 $\{(\mathrm{d}P)_p,(\mathrm{d}Q)_p\}$ 线性无关, 因而必要时缩小 $V$ 后, 可将 $(P,Q)$ 扩充为 $p$ 附近的一组局部坐标. 若 $m=2$, 则已得所需结论, 因为 $(Q,P)$ 是局部坐标, 且 $\{P,Q\}=X_P\cdot Q=1$.

若 $m>2$, 由于 $X_P$ 与 $T$ 不相切, 水平集 $P^{-1}(P(p))$ 与 $T$ 的交集为一个 $(m-2)$ 维流形 $S\subset T$. 由于 $\{P,Q\}=1$, 有 $X_Q\cdot P=\{Q,P\}=-1$, 所以 $X_Q$ 与 $S$ 不相切. 若 $q\in S$ 且 $\{v_1,\ldots,v_{m-2}\}$ 是 $T_qS$ 的一组基, 则 $\{(X_P)_q,(X_Q)_q,v_1,\ldots,v_{m-2}\}$ 是 $T_qM$ 的一组基. 此外,
\[
  \omega(X_P,v_i)=-\mathrm{d}P(v_i)=0
\]
(因为 $P$ 在 $S$ 上恒定), 同样有 $\omega(X_Q,v_i)=-\mathrm{d}Q(v_i)=0$. 由此得出矩阵 $\bigl(\omega(v_i,v_j)\bigr)$ 必非奇异, 即 $i^*\omega$ 必非退化, 其中 $i:S\to V$ 为包含映射. 由于 $\mathrm{d}i^*\omega=i^*\mathrm{d}\omega=0$, 可知 $(S,i^*\omega)$ 是辛流形. 给定任意 $F\in C^\infty(S)$, 可先令它沿 $X_Q$ 的流恒定而将其延拓到 $T$, 再令它沿 $X_P$ 的流恒定而将其延拓到 $V$. 由于 $[X_P,X_Q]=X_{\{P,Q\}}=X_1=0$, $X_P$ 与 $X_Q$ 的流彼此交换, 故也可以按相反的顺序延拓. 因此, 延拓后的函数 (仍记为 $F$) 满足 $\{P,F\}=X_P\cdot F=0$ 和 $\{Q,F\}=X_Q\cdot F=0$, 即 $X_F\cdot P=\{F,P\}=0$ 和 $X_F\cdot Q=\{F,Q\}=0$. 这意味着 $X_F$ 与 $S$ 相切, 所以 $X_F$ 在 $S$ 上与 $F$ 在 $(S,i^*\omega)$ 上确定的 Hamilton 向量场一致. 类似地, 两个函数 $F,G\in C^\infty(S)$ 延拓到 $V$ 后的 Poisson 括号 $\{F,G\}$ 满足
\begin{align*}
  X_P\cdot\{F,G\}
  &=\{P,\{F,G\}\}=\{\{P,F\},G\}+\{F,\{P,G\}\}\\
  &=\{0,G\}+\{F,0\}=0,
\end{align*}
同样有 $X_Q\cdot\{F,G\}=0$, 因而 $\{F,G\}$ 是 $(S,i^*\omega)$ 上的 Poisson 括号的延拓. 因此, 若 Darboux 定理对 $(S,i^*\omega)$ 成立, 即 $m-2=2n-2$, 且存在局部坐标 $(x^1,\ldots,x^{n-1},p_1,\ldots,p_{n-1})$, 使对 $i,j=1,\ldots,n-1$ 有 $\{x^i,x^j\}=\{p_i,p_j\}=0$ 和 $\{p_i,x^j\}=\delta_{ij}$, 则取 $x^n=Q$ 及 $p_n=P$, 就得到 $(M,\omega)$ 的结论.
\end{RiemannProof}
\RGfigure[.48\textwidth]{darboux-transverse-hypersurface}{Darboux 定理的证明}

事实上, 要有 Hamilton 流, 只需存在 Poisson 括号. 这提示我们可以作进一步推广.

\begin{RiemannDefinition}\label{def:poisson-manifold}
\textbf{Poisson 流形}是一个对 $(M,\{\cdot,\cdot\})$, 其中 $M$ 是可微流形, $\{\cdot,\cdot\}$ 称为 \textbf{Poisson 括号}, 是 $C^\infty(M)$ 上满足 \textbf{Leibniz 法则}的 Lie 括号, 即对任意 $\alpha,\beta\in\mathbb{R}$ 和 $F,G,H\in C^\infty(M)$ 满足:
\begin{enumerate}[label=(\roman*),leftmargin=*]
\item $\{F,G\}=-\{G,F\}$;
\item $\{\alpha F+\beta G,H\}=\alpha\{F,H\}+\beta\{F,H\}$;
\item $\{F,\{G,H\}\}+\{G,\{H,F\}\}+\{H,\{F,G\}\}=0$;
\item $\{F,GH\}=\{F,G\}H+\{F,H\}G$.
\end{enumerate}
\end{RiemannDefinition}

\begin{RiemannExample}\label{ex:basic-examples-of-poisson-manifolds}
\begin{enumerate}[label=(\arabic*),ref=\arabic*,leftmargin=*]
\item 任意辛流形 $(M,\omega)$ 自然是 Poisson 流形 $(M,\{\cdot,\cdot\})$ [参见 Exercise~\ref{ex:hamiltonian-mechanics}(\ref{ex:leibniz-rule-for-poisson-bracket})].
\item 任意光滑流形 $M$ 都可以赋予一个 Poisson 结构, 即平凡 Poisson 括号 $\{\cdot,\cdot\}:=0$. 对于辛结构, 即使 $M$ 的维数为偶数, 相应的结论也不成立 [参见 Exercise~\ref{ex:symmetry-and-reduction}(\ref{ex:canonical-brackets-and-symplectic-coordinates})].
\item 若 $\langle\cdot,\cdot\rangle$ 是 $\mathbb{R}^3$ 上的 Euclidean 内积, 则公式
\[
  \{F,G\}(x):=\langle x,\operatorname{grad}F\times\operatorname{grad}G\rangle
\]
在 $\mathbb{R}^3$ 上定义了一个非平凡的 Poisson 括号 (参见 Example~\ref{ex:rigid-body-as-reduced-hamiltonian-system}).
\end{enumerate}
\end{RiemannExample}

Poisson 括号的双线性及 Leibniz 法则意味着, 对任意 $F\in C^\infty(M)$, $\{F,\cdot\}$ 是一个导子 (因而是一个向量场). 这使我们可以定义 Hamilton 流.

\begin{RiemannDefinition}\label{def:hamiltonian-vector-field-on-poisson-manifold}
若 $(M,\{\cdot,\cdot\})$ 是 Poisson 流形, 且 $F\in C^\infty(M)$, 则由 $F$ 生成的 \textbf{Hamilton 向量场}是对任意函数 $G\in C^\infty(M)$ 满足
\[
  X_F\cdot G=\{F,G\}
\]
的向量场 $X_F\in\mathfrak{X}(M)$.
\end{RiemannDefinition}

\begin{RiemannProposition}\label{prop:poisson-and-vector-field-lie-homomorphism}
映射 $C^\infty(M)\ni F\mapsto X_F\in\mathfrak{X}(M)$ 是 $(C^\infty(M),\{\cdot,\cdot\})$ 与 $(\mathfrak{X}(M),[\cdot,\cdot])$ 之间的 Lie 代数同态, 即对所有 $\alpha,\beta\in\mathbb{R}$ 和 $F,G,H\in C^\infty(M)$:
\begin{enumerate}[label=(\roman*),leftmargin=*]
\item $X_{\alpha F+\beta G}=\alpha X_F+\beta X_G$;
\item $X_{\{F,G\}}=[X_F,X_G]$.
\end{enumerate}
\end{RiemannProposition}

\begin{RiemannProof}
性质 (i) 直接来自 Poisson 括号的双线性. 性质 (ii) 来自 Jacobi 恒等式, 因为对任意 $F,G,H\in C^\infty(M)$, 有
\begin{align*}
  X_{\{F,G\}}\cdot H
  &=\{\{F,G\},H\}=\{F,\{G,H\}\}-\{G,\{F,H\}\}\\
  &=X_F\cdot(X_G\cdot H)-X_G\cdot(X_F\cdot H)\\
  &=[X_F,X_G]\cdot H.
\end{align*}
\end{RiemannProof}

同态 $F\mapsto X_F$ 的核中的函数称为 \textbf{Casimir 函数}, 正是那些与其他所有函数都在 Poisson 括号下对合的函数 $F\in C^\infty(M)$, 即对所有 $G\in C^\infty(M)$ 满足 $\{F,G\}=0$ 的函数. 注意, Casimir 函数沿任意 Hamilton 流恒定. 同态 $F\mapsto X_F$ 的像是 Hamilton 向量场的集合, 特别地, 它构成 $(\mathfrak{X}(M),[\cdot,\cdot])$ 的一个 Lie 子代数.

\begin{RiemannExample}\label{ex:casimir-functions-and-spherical-level-sets}
\begin{enumerate}[label=(\arabic*),ref=\arabic*,leftmargin=*]
\item 若 $(M,\omega)$ 是辛流形, 则 Casimir 函数正是 (局部) 常函数.
\item 若 $\{\cdot,\cdot\}:=0$ 是光滑流形 $M$ 上的平凡 Poisson 括号, 则任意函数都是 Casimir 函数, 唯一的 Hamilton 向量场是零向量场.
\item 若 $\{\cdot,\cdot\}$ 是 $\mathbb{R}^3$ 上由公式
\[
  \{F,G\}(x):=\langle x,\operatorname{grad}F\times\operatorname{grad}G\rangle
\]
定义的 Poisson 括号, 则 $C(x):=\|x\|^2$ 是 Casimir 函数, 因为 $\operatorname{grad}C=2x$, 因而对任意光滑函数 $F\in C^\infty(\mathbb{R}^3)$ 有
\[
  \{C,F\}(x)=2\langle x,x\times\operatorname{grad}F\rangle
  =2\langle\operatorname{grad}F,x\times x\rangle=0.
\]
由此可知, Hamilton 向量场必与 $C$ 取常值的球面相切 (特别地, 必在原点为零).
\end{enumerate}
\end{RiemannExample}

由于 Poisson 括号可写为
\[
  \{F,G\}=X_F\cdot G=\mathrm{d}G(X_F)=-\mathrm{d}F(X_G),
\]
可见 $\{F,G\}(p)$ 分别是 $(\mathrm{d}F)_p$ 与 $(\mathrm{d}G)_p$ 的线性函数. 因此, Poisson 括号对所有 $p\in M$ 确定一个双线性映射 $B_p:T_p^*M\times T_p^*M\to\mathbb{R}$.

\begin{RiemannDefinition}\label{def:poisson-bivector}
满足
\[
  \{F,G\}=B(\mathrm{d}F,\mathrm{d}G)
\]
的反对称 $(0,2)$-张量场 $B$ 称为 \textbf{Poisson 双向量}.
\end{RiemannDefinition}

利用等同 $T_pM\cong(T_p^*M)^*$, 有
\begin{align*}
  X_F(\mathrm{d}G)&=\mathrm{d}G(X_F)=X_F\cdot G=\{F,G\}\\
                  &=B(\mathrm{d}F,\mathrm{d}G)=(\iota(\mathrm{d}F)B)(\mathrm{d}G),
\end{align*}
其中余向量与 Poisson 双向量的缩并, 按向量与交替张量的缩并同样定义 [参见第~\ref{chap:differential-forms}~章的 Exercise~\ref{ex:tensors}(\ref{ex:contraction-of-alternating-tensors})]. 因此有
\[
  X_F=\iota(\mathrm{d}F)B,
\]
所以在给定点 $p\in M$, Hamilton 向量场的所有可能取值所成的集合, 恰为映射 $T_p^*M\ni\omega\mapsto\iota(\omega)B\in T_pM$ 的像.

\begin{RiemannTheorem}[Kirillov]\label{thm:kirillov-symplectic-leaves-of-poisson-manifold}
设 $(M,\{\cdot,\cdot\})$ 是 Poisson 流形, 且映射 $T_p^*M\ni\omega\mapsto\iota(\omega)B\in T_pM$ 的秩为常数 (作为 $p\in M$ 的函数). 则 $M$ 具有由辛子流形 $(S,\omega_S)$ (称为\textbf{辛叶}) 构成的叶状结构, 使得
\[
  \{F,G\}(p)=\{F|_S,G|_S\}(p)
\]
对所有 $p\in M$ 成立, 其中 $S$ 是包含 $p$ 的叶.
\end{RiemannTheorem}

\begin{RiemannProof}
由于映射 $T_p^*M\ni\omega\mapsto\iota(\omega)B\in T_pM$ 的秩 $r$ 为常数, 其像 $\Sigma_p$ 对所有 $p\in M$ 都为 $r$ 维, 因而在 $M$ 上确定一个 $r$ 维分布 $\Sigma$. 按照构造, 所有 Hamilton 向量场都与该分布相容, 且显然对每个 $p\in M$, 都存在 $F_1,\ldots,F_r\in C^\infty(M)$, 使 $\Sigma$ 在 $p$ 的某个邻域内由 $X_{F_1},\ldots,X_{F_r}$ 张成. 由于对 $i,j=1,\ldots,r$ 都有 $[X_{F_i},X_{F_j}]=X_{\{F_i,F_j\}}$, 分布 $\Sigma$ 可积, 因而 $M$ 具有由 $r$ 维叶 $S$ 构成的叶状结构, 并且对所有 $p\in S$ 有 $T_pS=\Sigma_p$. 若 $\omega,\eta\in T_p^*M$, 则 $B(\omega,\eta)=\eta(\iota(\omega)B)=-\omega(\iota(\eta)B)$ 只依赖于 $\eta$ 与 $\omega$ 在 $\Sigma_p$ 上的限制, 即 $B$ 限制到 $\Sigma_p^*\times\Sigma_p^*$. 此外, 这一限制非退化, 因为映射 $\Sigma_p^*\ni\eta\mapsto\iota(\eta)B\in\Sigma_p$ 为满射. 接下来容易验证, 在每片叶 $S$ 上由 $B$ 在 $T^*S\times T^*S$ 上的限制所确定的 Poisson 括号, 来自 $S$ 上的一个辛形式 [参见 Exercise~\ref{ex:symmetry-and-reduction}(\ref{ex:symplectic-form-from-nondegenerate-poisson-bivector})].
\end{RiemannProof}

\begin{RiemannRemark}\label{rem:singular-symplectic-foliation}
Kirillov 定理在一般情形下仍成立, 此时映射 $T_p^*M\ni\omega\mapsto\iota(\omega)B\in T_pM$ 的秩不必为常数. 在此情形, 各辛叶不一定具有相同维数, 它们形成所谓的\textbf{奇异叶状结构}.
\end{RiemannRemark}

\begin{RiemannExample}\label{ex:symplectic-leaves-for-basic-poisson-structures}
\begin{enumerate}[label=(\arabic*),ref=\arabic*,leftmargin=*]
\item 若 $(M,\omega)$ 是辛流形, 则只有一片辛叶 ($M$ 本身).
\item 若 $\{\cdot,\cdot\}:=0$ 是光滑流形 $M$ 上的平凡 Poisson 括号, 则 Poisson 双向量恒为零, 辛叶是零维的点.
\item 若 $\{\cdot,\cdot\}$ 是 $\mathbb{R}^3$ 上由公式
\[
  \{F,G\}(x):=\langle x,\operatorname{grad}F\times\operatorname{grad}G\rangle
\]
定义的 Poisson 括号, 则 $x\in\mathbb{R}^3$ 处的 Poisson 双向量由
\[
  B(v,w)=\langle x,v\times w\rangle=\langle w,x\times v\rangle
\]
对任意 $v,w\in\mathbb{R}^3$ 给出, 其中利用 Euclidean 内积 $\langle\cdot,\cdot\rangle$ 作等同 $(\mathbb{R}^3)^*\cong\mathbb{R}^3$. 因此在 $x\in\mathbb{R}^3$ 处有
\[
  \iota(v)B=x\times v,
\]
所以 $B$ 在 $x$ 处的像是以原点为中心, 半径为 $\|x\|$ 的球面 $S_x$ 的切空间. 因而辛叶是球面 $S_x$ (包括原点, 它是一片奇异叶); 若 $x\ne0$, 则 $S_x$ 上的辛形式由
\[
  \omega(v,w)=\frac{1}{\|x\|^2}\langle x,v\times w\rangle
\]
对 $v,w\in T_xS_x$ 给出 (即 $\omega$ 是标准体积形式的 $\frac{1}{\|x\|}$ 倍). 事实上, 若 $F\in C^\infty(\mathbb{R}^3)$ 且 $v\in T_xS_x$, 则有
\begin{align*}
  \omega(X_F,v)
  &=\frac{1}{\|x\|^2}\langle x,X_F\times v\rangle
   =\frac{1}{\|x\|^2}\langle v,x\times(\iota(\operatorname{grad}F)B)\rangle\\
  &=\frac{1}{\|x\|^2}\langle v,x\times(x\times\operatorname{grad}F)\rangle\\
  &=\frac{1}{\|x\|^2}\bigl\langle v,\langle x,\operatorname{grad}F\rangle x
                              -\|x\|^2\operatorname{grad}F\bigr\rangle\\
  &=-\langle v,\operatorname{grad}F\rangle=-\mathrm{d}F(v).
\end{align*}
\end{enumerate}
\end{RiemannExample}

接下来考虑 Hamilton 流, 即由 Hamilton 向量场生成的流的几何性质. 与辛情形一样, 也有 Hamilton 形式的能量守恒.

\begin{RiemannProposition}\label{prop:poisson-hamiltonian-flow-preserves-generator}
Hamilton 流保持其生成函数.
\end{RiemannProposition}

\begin{RiemannProof}
若 $F\in C^\infty(M)$, 则
\[
  X_F\cdot F=\{F,F\}=-\{F,F\}=0.
\]
\end{RiemannProof}

回顾, 在辛情形下, Hamilton 流保持辛形式. 为得到这一性质在 Poisson 几何中的对应结论, 作如下定义.

\begin{RiemannDefinition}\label{def:poisson-map}
两个 Poisson 流形 $(M,\{\cdot,\cdot\})$ 和 $(N,\{\cdot,\cdot\})$ 之间的 \textbf{Poisson 映射} $f:M\to N$ 是可微映射, 且对所有 $F,G\in C^\infty(N)$ 满足
\[
  \{F,G\}\circ f=\{F\circ f,G\circ f\}.
\]
\end{RiemannDefinition}

如预期的那样, Poisson 映射保持 Hamilton 流.

\begin{RiemannProposition}\label{prop:poisson-map-intertwines-hamiltonian-fields}
若 $(M,\{\cdot,\cdot\})$ 和 $(N,\{\cdot,\cdot\})$ 为 Poisson 流形, $f:M\to N$ 为 Poisson 映射, 且 $F\in C^\infty(N)$, 则
\[
  f_*X_{F\circ f}=X_F.
\]
\end{RiemannProposition}

\begin{RiemannProof}
只需注意, 给定 $G\in C^\infty(N)$, 有
\[
  (f_*X_{F\circ f})\cdot G=X_{F\circ f}\cdot(G\circ f)
   =\{F\circ f,G\circ f\}=\{F,G\}=X_F\cdot G.
\]
\end{RiemannProof}

最后, 证明 Hamilton 流保持 Poisson 括号.

\begin{RiemannProposition}\label{prop:hamiltonian-flows-are-poisson-maps}
Hamilton 流是 Poisson 映射.
\end{RiemannProposition}

\begin{RiemannProof}
设 $\psi_t:M\to M$ 为函数 $F\in C^\infty(M)$ 生成的 Hamilton 流. 若 $G\in C^\infty(M)$ 为另一函数, 则有
\begin{align*}
  \frac{\mathrm{d}}{\mathrm{d}t}(G\circ\psi_t)
  &=\frac{\mathrm{d}}{\mathrm{d}t}(\psi_t^*G)
    =\left.\frac{\mathrm{d}}{\mathrm{d}s}\right|_{s=0}((\psi_{t+s})^*G)
    =\left.\frac{\mathrm{d}}{\mathrm{d}s}\right|_{s=0}((\psi_t\circ\psi_s)^*G)\\
  &=\left.\frac{\mathrm{d}}{\mathrm{d}s}\right|_{s=0}(\psi_s^*(\psi_t^*G))
    =X_F\cdot(\psi_t^*G)=\{F,\psi_t^*G\}.
\end{align*}
给定 $G,H\in C^\infty(M)$, 令 $K_t\in C^\infty(M)$ 为函数
\[
  K_t:=\{G,H\}\circ\psi_t-\{G\circ\psi_t,H\circ\psi_t\}
      =\psi_t^*\{G,H\}-\{\psi_t^*G,\psi_t^*H\}.
\]
显然 $K_0=0$. 由于 Poisson 括号双线性, 有
\begin{align*}
  \frac{\mathrm{d}}{\mathrm{d}t}K_t
  &=\frac{\mathrm{d}}{\mathrm{d}t}(\psi_t^*\{G,H\})
    -\left\{\frac{\mathrm{d}}{\mathrm{d}t}(\psi_t^*G),\psi_t^*H\right\}
    -\left\{\psi_t^*G,\frac{\mathrm{d}}{\mathrm{d}t}(\psi_t^*H)\right\}\\
  &=X_F\cdot(\psi_t^*\{G,H\})-\{\{F,\psi_t^*G\},\psi_t^*H\}
    -\{\psi_t^*G,\{F,\psi_t^*H\}\}\\
  &=X_F\cdot(\psi_t^*\{G,H\})-\{F,\{\psi_t^*G,\psi_t^*H\}\}
    =X_F\cdot K_t.
\end{align*}
将 $K_t$ 视为 $I\times M$ 上定义的函数 $K$, 其中 $I\subset\mathbb{R}$ 是 $\psi_t$ 的定义区间, 可见它满足
\[
  \begin{cases}
     \left(\dfrac{\partial}{\partial t}-X_F\right)\cdot K=0\\
     K(0,p)=0 & \text{对所有 $p\in M$}
  \end{cases}.
\]
从 $\{0\}\times M$ 出发沿 $\frac{\partial}{\partial t}-X_F$ 的积分曲线积分, 得到对所有 $t\in I$ 都有 $K_t=0$.
\end{RiemannProof}

现在可以讨论对称性与约化了.

\begin{RiemannDefinition}\label{def:poisson-and-hamiltonian-actions}
设 $G$ 为作用在 Poisson 流形 $(M,\{\cdot,\cdot\})$ 上的 Lie 群.
\begin{enumerate}[label=(\arabic*),leftmargin=*]
\item 若对每个 $g\in G$, 由 $A_g(p):=g\cdot p$ 给出的映射 $A_g:M\to M$ 是 Poisson 映射, 则称该作用是 \textbf{Poisson 作用};
\item 若对每个 $V\in\mathfrak{g}$, 都存在函数 $J(V)\in C^\infty(M)$, 使无穷小作用 $X^V$ 是由 $J(V)$ 生成的 Hamilton 向量场, 即 $X^V=X_{J(V)}$, 则称该作用是 \textbf{Hamilton 作用}.
\end{enumerate}
\end{RiemannDefinition}

若 $G$ 连通, 则 Proposition~\ref{prop:hamiltonian-flows-are-poisson-maps} 保证 Hamilton 作用是 Poisson 作用 [参见 Exercise~\ref{ex:symmetry-and-reduction}(\ref{ex:connected-hamiltonian-action-is-poisson})]. 注意, 由于 $X^V$ 是 $V$ 的线性函数, 可取 $J(V)$ 为 $V$ 的线性函数, 从而可将 $J$ 视为映射 $J:M\to\mathfrak{g}^*$. 该映射称为这一作用的\textbf{动量映射}.

\begin{RiemannTheorem}[Noether, Hamilton 形式]\label{thm:noether-for-hamiltonian-action}
若 Lie 群 $G$ 在 Poisson 流形 $(M,\{\cdot,\cdot\})$ 上的作用是具有动量映射 $J:M\to\mathfrak{g}^*$ 的 Hamilton 作用, 且 $H\in C^\infty(M)$ 是 $G$-不变的, 则 $J$ 沿 $H$ 的 Hamilton 流保持恒定.
\end{RiemannTheorem}

\begin{RiemannProof}
由于 $H$ 是 $G$-不变的, 对任意 $V\in\mathfrak{g}$ 有
\[
  X^V\cdot H=0\iff X_{J(V)}\cdot H=0\iff\{J(V),H\}=0\iff X_H\cdot J(V)=0.
\]
\end{RiemannProof}

\begin{RiemannExample}\label{ex:momentum-map-for-cotangent-lift}
下面这个重要的例子阐明了 Noether 定理的 Hamilton 形式与 Lagrange 形式之间的关系. 设 $M$ 为可微流形, $G$ 为作用在 $M$ 上的 Lie 群. 可通过下面的公式将这一作用提升到辛流形 (因而也是 Poisson 流形) $T^*M$ 上:
\[
  g\cdot\alpha=A_{g^{-1}}^*\alpha
\]
对所有 $\alpha\in T^*M$ 成立, 其中 $A_g:M\to M$ 由 $A_g(p)=g\cdot p$ 对所有 $p\in M$ 给出. 容易验证, 这个公式确实定义了 $G$ 在 $T^*M$ 上的作用, 将每个余切空间 $T_p^*M$ 映到 $T_{g\cdot p}^*M$.

设 $(x^1,\ldots,x^n)$ 为 $M$ 上的局部坐标, $(x^1,\ldots,x^n,p_1,\ldots,p_n)$ 为 $T^*M$ 上相应的局部坐标. 设 $(y^1,\ldots,y^m)$ 为 $G$ 上以单位元 $e\in G$ 为中心的局部坐标, 使 $(-y^1,\ldots,-y^m)$ 参数化由 $(y^1,\ldots,y^m)$ 参数化的元素的逆 (这容易用指数映射实现). 若在这些坐标下, $G$ 在 $M$ 上的作用 $A:G\times M\to M$ 由
\[
  \bigl(A^1(x^1,\ldots,x^n,y^1,\ldots,y^m),\ldots,
        A^n(x^1,\ldots,x^n,y^1,\ldots,y^m)\bigr)
\]
给出, 则有
\[
  A_{g^{-1}}^*\left(\sum_{i=1}^n p_i\,\mathrm{d}x^i\right)
  =\sum_{i,j=1}^n p_i\frac{\partial A^i}{\partial x^j}(x,-y)\,\mathrm{d}x^j,
\]
从而 $G$ 的作用到 $T^*M$ 的提升写为
\[
  \left(A^1(x,y),\ldots,A^n(x,y),
  \sum_{i=1}^n\frac{\partial A^i}{\partial x^1}(x,-y)p_i,\ldots,
  \sum_{i=1}^n\frac{\partial A^i}{\partial x^n}(x,-y)p_i\right).
\]
因此 $V:=\sum_{a=1}^m V^a\frac{\partial}{\partial y^a}$ 在 $T^*M$ 上的无穷小作用为
\begin{align*}
  &\sum_{i=1}^n X^i(x)\frac{\partial}{\partial x^i}
    -\sum_{i,j=1}^n\frac{\partial X^j}{\partial x^i}(x)p_j\frac{\partial}{\partial p_i}\\
  &\qquad=\sum_{i=1}^n\frac{\partial J}{\partial p_i}\frac{\partial}{\partial x^i}
    -\sum_{i=1}^n\frac{\partial J}{\partial x^i}\frac{\partial}{\partial p_i},
\end{align*}
其中
\[
  X^i(x)=\sum_{a=1}^m\frac{\partial A^i}{\partial y^a}(x,0)V^a
\]
是 $V$ 在 $M$ 上的无穷小作用的分量, 并且
\[
  J=\sum_{i=1}^n X^i(x)p_i.
\]
由此得出, $G$ 的作用到 $T^*M$ 的提升是 Hamilton 作用, 其动量映射 $J:T^*M\to\mathfrak{g}^*$ 由
\[
  J(\alpha)(V)=\alpha(X^V)
\]
给出, 其中 $X^V\in\mathfrak{X}(M)$ 是 $V$ 在 $M$ 上的无穷小作用. 注意, $J(V)$ 恰为 Noether 定理的 Lagrange 形式中的守恒量 $J^V$ 在 Legendre 变换下的像.
\end{RiemannExample}

\begin{RiemannTheorem}[Poisson 约化]\label{thm:poisson-reduction}
若 $G$ 在 $(M,\{\cdot,\cdot\})$ 上的作用是自由的, 适当的 Poisson 作用, 则 $M/G$ 自然是一个 Poisson 流形 (将 $C^\infty(M/G)$ 与 $C^\infty(M)$ 中的 $G$-不变函数等同), 且自然投影 $\pi:M\to M/G$ 为 Poisson 映射. 特别地, $\pi$ 将 $M$ 上 $G$-不变函数的 Hamilton 流映为 $M/G$ 上相应函数的 Hamilton 流.
\end{RiemannTheorem}

\begin{RiemannProof}
只需注意, 若该作用是 Poisson 作用, 则 $G$-不变函数的 Poisson 括号也是 $G$-不变的.
\end{RiemannProof}

若 $G$ 是 Lie 群, 则 $G$ 通过左乘在自身上作用, 且这一作用到 $T^*G$ 的提升是自由的, 适当的 Hamilton 作用. 如果 $G$ 还连通, 则该作用是 Poisson 作用, 由此得到如下结果.

\begin{RiemannTheorem}[Lie–Poisson 约化]\label{thm:lie-poisson-reduction}
若 $G$ 是连通 Lie 群, 则 $T^*G/G\cong\mathfrak{g}^*$ 上的商 Poisson 括号对所有 $F,H\in C^\infty(\mathfrak{g}^*)$ 由
\[
  \{F,H\}(\mu):=\mu([\mathrm{d}F,\mathrm{d}H])
\]
给出, 其中 $\mathrm{d}F,\mathrm{d}H\in\mathfrak{g}^{**}\cong\mathfrak{g}$. 若 $(p_1,\ldots,p_m)$ 是 $\mathfrak{g}^*$ 上与基 $\{\omega^1,\ldots,\omega^m\}$ 对应的线性坐标, 则
\[
  \{F,H\}=\sum_{a,b,c=1}^m p_aC_{bc}^a\frac{\partial F}{\partial p_b}\frac{\partial H}{\partial p_c},
\]
其中 $C_{bc}^a$ 是与 $\mathfrak{g}$ 的对偶基 $\{X_1,\ldots,X_m\}$ 相联系的结构常数.
\end{RiemannTheorem}

\begin{RiemannProof}
若将 $\{\omega^1,\ldots,\omega^m\}$ 看作 $G$ 上的左不变 1-形式, 则 $T^*G$ 上的典范辛势为
\[
  \theta=\sum_{a=1}^m p_a\omega^a
\]
(为简单起见, 将 $\omega^a$ 与 $\pi^*\omega^a$ 等同, 其中 $\pi:T^*G\to G$ 为自然投影). 由第~\ref{chap:curvature}~章的 Exercise~\ref{ex:cartan-structure-equations}(\ref{ex:structure-functions-of-frame}), 知道
\[
  \mathrm{d}\omega^a=-\frac{1}{2}\sum_{b,c=1}^m C_{bc}^a\omega^b\wedge\omega^c,
\]
因而典范辛形式为
\begin{align*}
  \omega=\mathrm{d}\theta
  &=\sum_a^m\mathrm{d}p_a\wedge\omega^a-\frac{1}{2}
    \sum_{a,b,c=1}^m p_aC_{bc}^a\omega^b\wedge\omega^c\\
  &=\sum_a^m\mathrm{d}p_a\otimes\omega^a-\sum_a^m\omega^a\otimes\mathrm{d}p_a
    -\sum_{a,b,c=1}^m p_aC_{bc}^a\omega^b\otimes\omega^c.
\end{align*}
若 $F\in C^\infty(T^*G)$ 是 $G$-不变的, 则它只依赖纤维方向的坐标 $(p_1,\ldots,p_m)$, 因此
\[
  \mathrm{d}F=\sum_{a=1}^m\frac{\partial F}{\partial p_a}\,\mathrm{d}p_a.
\]
令
\[
  X_F:=\sum_{a=1}^m\xi^aX_a+\sum_{a=1}^m\eta_a\frac{\partial}{\partial p_a},
\]
其中 $\{X_1,\ldots,X_m\}$ 为 $\mathfrak{g}$ 的对偶基, 则
\[
  \iota(X_F)\omega=-\sum_{a=1}^m\xi^a\,\mathrm{d}p_a
    -\sum_{a,b,c=1}^m p_aC_{bc}^a\xi^b\omega^c+\sum_{a=1}^m\eta_a\omega^a.
\]
由 $\iota(X_F)\omega=-\mathrm{d}F$, 得到
\[
  X_F=\sum_{a=1}^m\frac{\partial F}{\partial p_a}X_a
    +\sum_{a,b,c=1}^m p_aC_{bc}^a\frac{\partial F}{\partial p_b}\frac{\partial}{\partial p_c},
\]
所以若 $H\in C^\infty(T^*G)$ 也是 $G$-不变的, 则
\[
  \{F,H\}=X_F\cdot H=\sum_{a,b,c=1}^m p_aC_{bc}^a
      \frac{\partial F}{\partial p_b}\frac{\partial H}{\partial p_c}.
\]
注意, 作为 $\mathfrak{g}^*$ 上的余向量, $\mathrm{d}p_a=X_a$. 所以由结构函数 $C_{bc}^a$ 的定义,
\begin{align*}
  \{F,H\}
  &=\sum_{a,b,c=1}^m p_a\,\omega^a([X_b,X_c])
      \frac{\partial F}{\partial p_b}\frac{\partial H}{\partial p_c}\\
  &=\sum_{a,b,c=1}^m p_a\,\omega^a\left(
        \left[\frac{\partial F}{\partial p_b}X_b,\frac{\partial H}{\partial p_c}X_c\right]\right)\\
  &=\sum_a^m p_a\,\omega^a([\mathrm{d}F,\mathrm{d}H]).
\end{align*}
\end{RiemannProof}

\begin{RiemannExample}\label{ex:rigid-body-as-reduced-hamiltonian-system}
$T^*SO(3)$ 上的 Lie–Poisson 约化给出 $\mathfrak{so}(3)^*\cong(\mathbb{R}^3)^*\cong\mathbb{R}^3$ 上的 Poisson 括号
\[
  \{F,G\}(x):=\langle x,\nabla F\times\nabla G\rangle,
\]
其中利用 Lemma~\ref{lem:so-three-cross-product-isomorphism} 将 $\mathfrak{so}(3)$ 与 $(\mathbb{R}^3,\times)$ 等同, 并利用 Euclidean 内积 $\langle\cdot,\cdot\rangle$ 作等同 $(\mathbb{R}^3)^*\cong\mathbb{R}^3$.
\end{RiemannExample}

\begin{RiemannExercise}\label{ex:symmetry-and-reduction}
\begin{enumerate}[label=(\arabic*),ref=\arabic*,leftmargin=*]
\item\label{ex:height-hamiltonian-on-sphere}
考虑由通常的体积形式在
\[
  S^2=\{(x,y,z)\in\mathbb{R}^3\mid x^2+y^2+z^2=1\}
\]
上确定的辛结构. 求由函数 $H(x,y,z)=z$ 生成的 Hamilton 流.
\item\label{ex:canonical-brackets-and-symplectic-coordinates}
设 $(M,\omega)$ 为辛流形. 证明:
\begin{enumerate}[label=(\alph*),leftmargin=*]
\item $\omega=\sum_{i=1}^n\mathrm{d}p_i\wedge\mathrm{d}x^i$, 当且仅当对 $i,j=1,\ldots,n$ 有 $\{x^i,x^j\}=\{p_i,p_j\}=0$ 及 $\{p_i,x^j\}=\delta_{ij}$;
\item $M$ 可定向;
\item 若 $M$ 紧, 则 $\omega$ 不可能是恰当的. (\textbf{注}: 特别地, 若 $M$ 紧, 且 $M$ 上所有闭 2-形式都是恰当的, 则 $M$ 不容许辛结构; 所有满足 $n>1$ 的偶数维球面 $S^{2n}$ 都属于这种情形.)
\end{enumerate}
\item\label{ex:magnetic-term-in-symplectic-form}
设 $(M,\langle\cdot,\cdot\rangle)$ 为 Riemann 流形, $\alpha\in\Omega^1(M)$ 为 1-形式, $U\in C^\infty(M)$ 为可微函数.
\begin{enumerate}[label=(\alph*),leftmargin=*]
\item 证明 $\widetilde\omega:=\omega+\pi^*\mathrm{d}\alpha$ 是 $T^*M$ 上的一个辛形式, 其中 $\omega$ 为典范辛形式, $\pi:T^*M\to M$ 为自然投影 ($\widetilde\omega$ 称为\textbf{带磁项的典范辛形式}).
\item 证明函数 $\widetilde H\in C^\infty(T^*M)$ 关于辛形式 $\widetilde\omega$ 所生成的 Hamilton 流由以下方程给出:
\[
  \begin{cases}
    \dot x^i=\dfrac{\partial\widetilde H}{\partial p_i}\\[.7ex]
    \displaystyle\dot p_i=-\frac{\partial\widetilde H}{\partial x^i}
    +\sum_{j=1}^n\left(\frac{\partial\alpha_j}{\partial x^i}
                  -\frac{\partial\alpha_i}{\partial x^j}\right)\dot x^j
  \end{cases}.
\]
\item 由
\[
  F(\xi):=\xi-\alpha_p
\]
对 $\xi\in T_p^*M$ 给出的映射 $F:T^*M\to T^*M$ 是保纤维的微分同胚. 证明 $F$ 将 Exercise~\ref{ex:hamiltonian-mechanics}(\ref{ex:conservative-system-with-magnetic-term}) 中的 Hamilton 流映为 $\widetilde H$ 关于辛形式 $\widetilde\omega$ 的 Hamilton 流, 其中
\[
  \widetilde H(\xi):=\frac{1}{2}\langle\xi,\xi\rangle+U(p)
\]
对 $\xi\in T_p^*M$ 成立. (\textbf{注}: 由于两个流在 $M$ 上的投影一致, 可见可以通过改变 Lagrange 函数或辛形式来引入磁项.)
\end{enumerate}
\item\label{ex:symplectic-form-from-nondegenerate-poisson-bivector}
设 $(M,\{\cdot,\cdot\})$ 为 Poisson 流形, $B$ 为 Poisson 双向量, $(x^1,\ldots,x^n)$ 为 $M$ 上的局部坐标. 证明:
\begin{enumerate}[label=(\alph*),leftmargin=*]
\item $B$ 在这些局部坐标下可写为
\[
  B=\sum_{i,j=1}^n B^{ij}\frac{\partial}{\partial x^i}\otimes\frac{\partial}{\partial x^j},
\]
其中对 $i,j=1,\ldots,n$ 有 $B^{ij}=\{x^i,x^j\}$;
\item $F\in C^\infty(M)$ 生成的 Hamilton 向量场可写为
\[
  X_F=\sum_{i,j=1}^p B^{ij}\frac{\partial F}{\partial x^i}\frac{\partial}{\partial x^j};
\]
\item $B$ 的各分量必须满足
\[
  \sum_{l=1}^n\left(
    B^{il}\frac{\partial B^{jk}}{\partial x^l}
   +B^{jl}\frac{\partial B^{ki}}{\partial x^l}
   +B^{kl}\frac{\partial B^{ij}}{\partial x^l}\right)=0
\]
对所有 $i,j,k=1,\ldots,n$ 成立;
\item 若 $\{\cdot,\cdot\}$ 来自辛形式 $\omega$, 则 $(B^{ij})=-(\omega_{ij})^{-1}$;
\item 若 $B$ 非退化, 则它来自一个辛形式.
\end{enumerate}
\item\label{ex:action-angle-coordinates}
(\emph{作用--角坐标}) 设 $(M,\omega)$ 为辛流形, $F=(F_1,\ldots,F_n):M\to\mathbb{R}^n$ 为一组独立的对合第一积分, 且具有紧水平集 ($n$ 维环面). 选一个不变环面 $T_0$, 一点 $\alpha\in T_0$, 以及在 $\alpha$ 处与 $T_0$ 横截的 $n$ 维子流形 $N\subset M$ (即 $T_\alpha M=T_\alpha T\oplus T_\alpha N$). 在每个靠近 $T_0$ 的不变环面 $T$ 上, 由等同 $T\cong\mathbb{R}^n/\mathbb{Z}^n$ (来自与 $F$ 相关联的 $\mathbb{R}^n$ 作用) 确定坐标 $x=(x^1,\ldots,x^n)$, 令 $x=0$ 在 $N$ 上成立, 从而固定这些坐标. 这样, 就得到 $T$ 的一个邻域内的局部坐标 $(x,F)$. 证明:
\begin{enumerate}[label=(\alph*),leftmargin=*]
\item 在这些坐标下, Poisson 双向量的分量为
\[
  (B)=\left(\begin{array}{c|c}
    * & -A^t\\ \hline
    A & 0
  \end{array}\right),
\]
因而辛形式的分量为
\[
  (\omega)=\left(\begin{array}{c|c}
    0 & -A^{-1}\\ \hline
    (A^{-1})^t & *
  \end{array}\right),
\]
其中 $A=A(F)$ 是由下式定义的矩阵 $A=(a_{ij})$:
\[
  X_{F_i}=\sum_{j=1}^n a_{ij}\frac{\partial}{\partial x^j};
\]
\item 可选取新坐标 $J=J(F)$, 使 $\{J_i,x^j\}=\delta_{ij}$;
\item $\{x^i,x^j\}$ 只是 $J$ 的函数;
\item 可选取新坐标 $y=x+z(J)$, 使 $\{y^i,y^j\}=0$ 且 $\{J_i,y^j\}=\delta_{ij}$;
\item 在 $T_0$ 的一个邻域中存在 1-形式 $\theta$, 使 $\omega=\mathrm{d}\theta$, 且 $J_i=\oint_{\gamma_i}\theta$, 其中 $\gamma_i$ 是 $y^i$ 轴在各不变环面 $T\cong\mathbb{R}^n/\mathbb{Z}^n$ 上的投影.

(\textbf{提示}: 将用到 Poincaré 引理 --- 参见第~\ref{chap:differential-forms}~章的 Exercise~\ref{ex:differential-forms}(\ref{ex:homotopy-operator-and-poincare-lemma}).)
\end{enumerate}
\item\label{ex:connected-hamiltonian-action-is-poisson}
设 $G$ 为连通 Lie 群, $U\subset G$ 为单位元的邻域. 证明:
\begin{enumerate}[label=(\alph*),leftmargin=*]
\item $G=\bigcup_{n=1}^{+\infty}U^n$, 其中 $U^n=\{g_1\cdots g_n\mid g_1,\ldots,g_n\in U\}$;
\item 若 $G$ 在 Poisson 流形 $(M,\{\cdot,\cdot\})$ 上的作用是 Hamilton 作用, 则该作用是 Poisson 作用.
\end{enumerate}
\item\label{ex:relative-equilibrium-on-group-orbit}
设连通 Lie 群 $G$ 在 Poisson 流形 $(M,\{\cdot,\cdot\})$ 上有自由且适当的 Hamilton 作用, $H\in C^\infty(M)$ 为 $G$-不变的函数. 证明若 $p\in M/G$ 是 $\pi_*X_H\in\mathfrak{X}(M/G)$ 的不动点 (其中 $\pi:M\to M/G$ 为商映射), 则 $X_H$ 在 $\pi^{-1}(p)$ 上的流由 $G$ 的单参数子群的轨道给出.
\item\label{ex:poisson-reduction-for-central-force}
Lie 群 $SO(2)\cong S^1$ 以
\[
  e^{i\varphi}\cdot(r,\theta)=(r,\theta+\varphi)
\]
作用在 $M=\mathbb{R}^2\setminus\{(0,0)\}$ 上, 其中 $(r,\theta)$ 是 $M$ 上的极坐标, $\theta,\varphi$ 按模 $2\pi$ 理解.
\begin{enumerate}[label=(\alph*),leftmargin=*]
\item 写出 $V\in\mathfrak{so}(2)\cong\mathbb{R}$ 的无穷小作用 $X^V\in\mathfrak{X}(M)$ 的表达式.
\item 求这一作用到余切丛的提升的动量映射.
\item 在通常坐标 $(r,\theta,p_r,p_\theta)$ 下, 写出具有典范辛结构的 $T^*M$ 的 Poisson 双向量的表达式.
\item 求 Poisson 流形 $Q:=T^*M/SO(2)\cong\mathbb{R}^3$ 的 Poisson 双向量. 此流形的辛叶是什么? 举一个非恒定 Casimir 函数的例子.
\item 考虑由下式给出的 Hamilton 函数 $H:T^*Q\to\mathbb{R}$:
\[
  H(r,\theta,p_r,p_\theta)=\frac{p_r^2}{2}+\frac{p_\theta^2}{2r^2}+u(r).
\]
证明 $H$ 是 $SO(2)$-不变的, 并求它在约化 Poisson 流形 $Q$ 上的 Hamilton 流.
\item 利用 Noether 定理, 得到 $H$ 在 $T^*M$ 上的 Hamilton 流所保持的一个量.
\end{enumerate}
\item\label{ex:poisson-reduction-of-hyperbolic-geodesic-flow}
回顾, 上半平面 $H=\{(x,y)\in\mathbb{R}^2\mid y>0\}$ 具有 Lie 群结构, 其运算由
\[
  (x,y)\cdot(z,w):=(yz+x,yw)
\]
给出. 双曲平面对应于 $H$ 上的左不变度量
\[
  g:=\frac{1}{y^2}(\mathrm{d}x\otimes\mathrm{d}x+\mathrm{d}y\otimes\mathrm{d}y)
\]
[参见第~\ref{chap:differentiable-manifolds}~章的 Exercise~\ref{ex:lie-groups}(\ref{ex:affine-maps-lie-group}) 及第~\ref{chap:riemannian-manifolds}~章的 Exercise~\ref{ex:levi-civita-connection}(\ref{ex:hyperbolic-plane})]. 因此, 测地线由 Hamilton 函数 $K:T^*H\to\mathbb{R}$ 确定, 其中
\[
  K(x,y,p_x,p_y)=\frac{y^2}{2}(p_x^2+p_y^2).
\]
\begin{enumerate}[label=(\alph*),leftmargin=*]
\item 求 $H$ 在自身上的左平移作用到 $T^*H$ 的提升, 并核实它保持 Hamilton 函数 $K$.
\item 证明函数
\[
  F(x,y,p_x,p_y)=yp_x\qquad\text{和}\qquad G(x,y,p_x,p_y)=yp_y
\]
也是 $H$-不变的, 并利用这一点求 $T^*H/H$ 上的商 Poisson 结构. 这是辛流形吗?
\item 写出 $H$ 在 $T^*H$ 上作用的动量映射的表达式, 并利用它求测地线方程的一个非平凡第一积分 $I$. 证明满足 $K=E$, $p_x=l$ 和 $I=m$ 的测地线在 $H$ 上的投影满足方程
\[
  l^2x^2+l^2y^2-2lmx+m^2=2E.
\]
假定 $l\ne0$, 这些曲线是什么?
\end{enumerate}
\item\label{ex:lie-poisson-reduction-of-euler-top}
回顾, Euler 陀螺是由 Lagrange 函数 $L:TSO(3)\to\mathbb{R}$ 确定的力学系统, 其中
\[
  L=\frac{1}{2}\langle I\Omega,\Omega\rangle,
\]
$\Omega$ 是由下列等同得到的纤维上的左不变坐标:
\[
  T_SSO(3)=(\mathrm{d}L_S)(\mathfrak{so}(3))\cong\mathfrak{so}(3)\cong\mathbb{R}^3.
\]
\begin{enumerate}[label=(\alph*),leftmargin=*]
\item 证明若利用 Euclidean 内积 $\langle\cdot,\cdot\rangle$ 将 $(\mathbb{R}^3)^*$ 与 $\mathbb{R}^3$ 等同, 则 Legendre 变换写为
\[
  P=I\Omega,
\]
其中 $P$ 是 $T^*SO(3)$ 上相应的左不变坐标.
\item 写出约化 Poisson 流形 $T^*SO(3)/SO(3)\cong\mathbb{R}^3$ 上的 Hamilton 方程. 辛叶是什么? 举一个非恒定 Casimir 函数的例子.
\item 求 $SO(3)$ 在自身上的左平移作用到 $T^*SO(3)$ 的提升的动量映射.
\end{enumerate}
\item\label{ex:poisson-brackets-of-body-momentum-and-direction}
设 $(P^1,P^2,P^3)$ 是 $T^*SO(3)$ 的纤维上通常的左不变坐标, 并考虑对每个 $S\in SO(3)$ 由下式定义的函数 $(\Gamma^1,\Gamma^2,\Gamma^3)$:
\[
  \gamma=S\Gamma,
\]
其中 $\gamma\in\mathbb{R}^3$ 为固定向量. 证明对 $i,j=1,2,3$ 有:
\begin{enumerate}[label=(\alph*),leftmargin=*]
\item $\{P^i,P^j\}=\sum_{k=1}^3\varepsilon_{ijk}P^k$;
\item $\{\Gamma^i,\Gamma^j\}=0$;
\item $\{P^i,\Gamma^j\}=\sum_{k=1}^3\varepsilon_{ijk}\Gamma^k$,
\end{enumerate}
其中
\[
  \varepsilon_{ijk}=
  \begin{cases}
    +1 & \text{若 $(i,j,k)$ 是 $(1,2,3)$ 的偶排列}\\
    -1 & \text{若 $(i,j,k)$ 是 $(1,2,3)$ 的奇排列}\\
     0 & \text{其他情形}.
  \end{cases}
\]
(\textbf{提示}: 证明沿 Euler 陀螺的任意运动有 $\dot\Gamma=\Gamma\times\Omega$, 其中 $\Omega$ 为 Euler 陀螺的随体参考系中的角速度. 再把 $\frac{(P^i)^2}{2}$ 看作 Euler 陀螺的 Hamilton 函数在 $I_i=1$ 且对 $j\ne i$ 有 $I_j\to+\infty$ 时的极限.)
\item\label{ex:poisson-reduction-of-heavy-rigid-body}
若在 Exercise~\ref{ex:symmetry-and-reduction}(\ref{ex:poisson-brackets-of-body-momentum-and-direction}) 中取 $\gamma=ge_z$, 其中 $g$ 为 (恒定的) 重力加速度, 则质量为 $M$, 转动惯量为 $I$ 的有定点刚体, 若其质心在随体参考系中的位置向量为 $L\in\mathbb{R}^3$, 则它的运动由
\[
  H=\frac{1}{2}\langle P,I^{-1}P\rangle+M\langle\gamma,SL\rangle
\]
的 Hamilton 流给出.
\begin{enumerate}[label=(\alph*),leftmargin=*]
\item $S^1$ 在 $SO(3)$ 上的作用由 $e^{i\theta}\cdot S=R_\theta S$ 确定, 其中
\[
  R_\theta:=\begin{pmatrix}
    \cos\theta & -\sin\theta & 0\\
    \sin\theta & \cos\theta & 0\\
    0 & 0 & 1
  \end{pmatrix}
\]
(对应于绕 $z$ 轴的旋转). 证明对于这一作用到 $T^*SO(3)$ 的提升, $H$ 是 $S^1$-不变的.
\item 求这一作用的动量映射.
\item 证明函数 $P$ 和 $\Gamma$ 都是 $S^1$-不变的, 且商流形 $T^*SO(3)/S^1\cong(SO(3)/S^1)\times\mathbb{R}^3\cong S^2\times\mathbb{R}^3$ 上的 Poisson 括号由这些函数的 Poisson 括号确定.
\item 用函数 $P$ 和 $\Gamma$ 写出商流形上的运动方程, 并举一个非恒定 Casimir 函数的例子.
\end{enumerate}
\end{enumerate}
\end{RiemannExercise}
